% OB-SECTION {"id":"benchmark","label":"sec-benchmark","level":2,"order":3,"title":"Benchmark"}

\section{Benchmark}
\label{sec-benchmark}
% OB-EDIT-BEGIN openbarry/markdown@1 id=benchmark
Our benchmark evaluates language model diversity using two distinct datasets: \textbf{NB-Curated}, which contains 100 prompts manually curated by the authors, and \textbf{NB-WildChat}, which consists of 1,000 prompts automatically curated from real user interactions with ChatGPT.
LLM performance is evaluated based on the quality and novelty of the set of generations produced by the LLM.

\subsection{Dataset curation}

To test language models' capabilities to produce novel generations, we need a dataset of prompts that elicit diverse responses.

\paragraph{NB-Curated.}
The curated dataset is designed to contain four distinct categories where diversity is expected.

\begin{itemize}
\item \textbf{Randomness:} prompts that involve randomizing over a set of options. Example: \emph{Roll a make-believe 20-sided die.}
\item \textbf{Factual knowledge:} prompts that request underspecified factual information. Example: \emph{List a capital city in Africa.}
\item \textbf{Creative writing:} prompts that involve generating a creative form of text. Example: \emph{Tell me a riddle.}
\item \textbf{Subjectivity:} prompts that request subjective answers or opinions. Example: \emph{What's the best car to get in 2023?}
\end{itemize}

\obtodo[id=nb-curated-write, created=2024-10-14T16:03:12.000Z, updated=2024-11-22T21:47:03.000Z, completed=2024-11-22T21:47:03.000Z]{Write 100 NB-Curated prompts across randomness, factual knowledge, creative writing, and subjectivity.}
\obtodo[id=nb-human-annotate, created=2024-11-12T10:08:00.000Z, updated=2024-12-18T19:26:41.000Z, completed=2024-12-18T19:26:41.000Z]{Collect eight author responses for each of the 100 curated prompts.}
\obcomment[kind=result]{8 humans x 100 prompts. Homogeneous annotator pool (same institution), so treat this as a lower bound on human diversity, not a population estimate.}

The authors of the paper then write responses to each of the prompts, resulting in 8 human responses per prompt.
Since this is a rather homogeneous set of human annotators (all work/study at the same institution), we expect this response set to represent an approximate lower bound on actual human diversity.
However, as we will show in Section~\ref{sec-evaluation}, even this group produced more diverse responses than most state-of-the-art language models.

\paragraph{NB-WildChat.}
WildChat is a collection of 1M real user conversations with ChatGPT~\cite{zhao2023wildchat}.
We cannot use this dataset directly because the vast majority of its prompts do not allow for multiple valid answers.
During processing, we first used Llama Guard 3~\cite{chi2024} to filter out inappropriate user prompts.
A small number of IPs contributed a disproportionately high number of interactions, leading to topic bias (e.g., towards anime); we deduplicated the dataset based on user IP.
After preprocessing, we used GPT-4o as a classifier to select prompts expected to elicit diverse generations (prompt in Appendix~\ref{sec-appendix}).
To make sure that the selected prompts match the desired criteria, we manually labeled 100 prompts from the dataset (with 50 accepted and 50 rejected by GPT-4o), and found an 85\% agreement between the classifier and the human labels.

\obtodo[id=nb-wildchat-filter, created=2024-12-02T11:20:00.000Z, updated=2025-01-16T08:55:17.000Z, completed=2025-01-16T08:55:17.000Z]{Filter WildChat-1M with Llama Guard 3, IP dedup, then GPT-4o diversity classifier to 1,000 prompts.}
\obtodo[id=nb-wildchat-verify, created=2025-01-16T09:12:00.000Z, updated=2025-01-24T17:03:09.000Z, completed=2025-01-24T17:03:09.000Z]{Manually label 100 WildChat candidates, 50 accepted and 50 rejected by GPT-4o.}
\obcomment[kind=result]{85 percent agreement with GPT-4o on the 100-prompt audit. Good enough to ship the 1,000-prompt split.}

\subsection{Evaluating novelty}

Models with diversity alone would not be sufficient for users: each of the individual generations should also be of high quality, with the ultimate goal of maximizing the cumulative utility of a user who interacts with the model.

\paragraph{Partitioning the output space.}
Traditional metrics like $n$-gram overlap or embedding similarity often fail to capture meaningful diversity.
We mainly consider \emph{functional equivalence}: two generations are different if and only if a user who has seen one generation would likely benefit from seeing the other.
For example, in a story generation task, if two generations are retelling the original story with identical plot but different character names, then we consider them to be functionally equivalent.

The authors annotated 1,100 pairs of generations conditioned on prompts from both NB-Curated and NB-WildChat, sampled from a diverse set of models.
From these annotated pairs, we used 1,000 for training and fine-tuned a \texttt{deberta-v3-large} model~\cite{he2021deberta} to predict binary functional equivalence.
On a held-out test set of 100 pairs, the classifier achieved 79\% accuracy compared to human labels, with an F1 score of 0.811.

\obtodo[id=nb-equiv-annotate, created=2025-01-06T14:00:00.000Z, updated=2025-01-31T22:18:05.000Z, completed=2025-01-31T22:18:05.000Z]{Annotate 1,100 generation pairs for functional equivalence. Train on 1,000, hold out 100.}
\obtodo[id=nb-deberta-train, created=2025-02-01T09:40:00.000Z, updated=2025-02-11T13:27:51.000Z, completed=2025-02-11T13:27:51.000Z]{Fine-tune deberta-v3-large on the 1,000 training pairs and pick an operating threshold.}
\obcomment[kind=result]{Held-out: 79 percent accuracy, F1 0.811. BLEU/ROUGE/BERTScore all worse; keep DeBERTa as the partitioner. Classifier truncates each response to 128 tokens --- note that in the repo README.}

With the equivalence classifier, we partition $k$ sample generations into classes.
For each new generation, we compare it against a random generation from each existing class.
If the classifier determines it is functionally equivalent to any existing class, we assign it to the first such class found; otherwise, we create a new class.
We define
\begin{equation}
\mathrm{distinct}_k := \bigl|\bigl\{c_i \mid i \in [k]\bigr\}\bigr|,
\end{equation}
where $c_i$ is the equivalence class of the $i$-th generation.

\paragraph{Unified measure of diversity and quality.}
Inspired by rank-biased precision~\cite{moffat2008}, we assume the user has patience $p \in [0,1]$: after each generation they request another with probability $p$.
The utility of a generation that is functionally equivalent to a previous generation is zero:
\begin{equation}
\mathrm{utility}_k := \frac{1-p}{1-p^k} \sum_{i=1}^{k} p^{i-1} \cdot \mathbf{1}[c_i \neq c_j, \forall j < i] \cdot u_i.
\end{equation}
In our evaluation setup, we set the patience level to $0.8$.
As $p \to 0$, the metric converges to quality evaluation of a single generation.

\paragraph{Computing utility.}
We use the Skywork-Reward-Gemma-2-27B-v0.2 model~\cite{liu2024skywork} to score the quality of individual generations, chosen as a top RewardBench~\cite{lambert2024} model.
Reward scores are mapped into $\{1,2,\ldots,10\}$ by calibrating against GPT-4 quality scores on 2,400 MT-Bench~\cite{zheng2023mtbench} generations.

\obtodo[id=nb-utility-calibrate, created=2025-02-12T10:05:00.000Z, updated=2025-02-20T18:44:12.000Z, completed=2025-02-20T18:44:12.000Z]{Calibrate Skywork-Reward scores to a 1-10 utility using MT-Bench GPT-4 judgments.}
% OB-EDIT-END id=benchmark
